\section{Introduction}
Quantum machine learning (QML) explores whether quantum-state representations, parameterized unitary transformations, and quantum measurement statistics can provide expressive and trainable models for learning tasks. In a variational quantum classifier (VQC) or quantum neural network (QNN), classical input features are embedded into the state space of a register of qubits, parameterized quantum logic gates transform the state, and observables are measured to produce prediction logits. The entire circuit is typically trained within a hybrid quantum-classical loop using gradient descent [1,2].

Despite considerable enthusiasm, empirical investigations frequently reveal that classical neural networks consistently outperform variational quantum models on standard classical data benchmarks [3]. A high accuracy score on a simplistic dataset cannot establish quantum utility. A meaningful comparison demands matched data partitions, identical training/test budgets, paired statistical tests, capacity-controlled baselines, and a clear account of inductive biases and computational work.

\subsection{Research questions}
This case study addresses four concrete research questions:
\begin{enumerate}
  \item How does a variational quantum classifier compare with classical MLPs across three distinct data regimes: quantum-native phase classification (PennyLane TFIM, Qiskit Phase of Matter), high-dimensional non-linear quantum kernels (Havl\'i\v{c}ek ad-hoc benchmark), and classical small-data benchmarks (Breast Cancer, downsampled MNIST)?
  \item Why do classical neural networks overwhelmingly outperform QNNs on generic classical tabular data?
  \item Under what geometric data structures or circuit configurations do QNNs exhibit competitive inductive bias advantages over capacity-constrained classical networks?
  \item What changes when circuit depth is scaled, when entangling gates are removed, or when exact analytic expectations are replaced by finite-shot measurement sampling?
\end{enumerate}

\subsection{What this comparison can and cannot show}
The unit of comparison in this paper is a trained classifier under a specified data split, feature representation, optimization routine, and simulator. It is not a comparison between a physical quantum processor and a CPU. The quantum circuit is simulated classically, and both models run within a classical training loop. This experiment can answer practical questions about decision boundaries, optimization, circuit design, measurement sampling, and inductive bias. It cannot show that quantum hardware reduces computational wall-clock time or computational complexity on classical inputs.

The main practical design choice is to keep the comparison controlled: all models receive identical standardized continuous features ($d=2$ for 2-qubit kernel benchmarks, $d=4$ for 4-qubit quantum state and PCA-projected small-data benchmarks). The QNN encodes those features as $R_y$ rotation angles; the MLP receives the standardized numeric values. Preprocessing is fit only on training data, and model comparisons reuse the same held-out examples for every seed.

\subsection{Contributions of this case study}
The study makes five bounded contributions:
\begin{enumerate}
  \item Evaluates five contrasting, domain-purposeful benchmarks spanning quantum state classification, non-linear quantum kernels, and classical small-data regimes: PennyLane TFIM, Qiskit Phase of Matter, Havl\'i\v{c}ek ad-hoc kernel, Breast Cancer Wisconsin, and downsampled MNIST digits (0 vs.\ 1).
  \item Precedes all modeling with automated Exploratory Data Analysis (EDA) quantifying column schemas, data types, target designations, class balance, geometric separation metrics, inter-feature correlations, and PCA spectral dimensionality.
  \item Employs capacity-controlled classical networks (9 parameters for 2D benchmarks, 13 parameters for 4D benchmarks) alongside overparameterized wide MLPs (193 and 225 parameters) to match the parameter budgets of depth-2 QNNs (7 and 13 parameters).
  \item Isolates circuit depth scaling (depths 1--3) and entanglement ablations (with vs. without CNOT) across all five datasets.
  \item Formulates the theoretical and empirical principles governing why classical MLPs dominate classical data while detailing the foundational frontiers where quantum models hold empirical parity or potential advantage.
\end{enumerate}
